\title{Direct Preference Optimization: Your Language Model is Secretly a Reward Model}

\begin{abstract}
Although large-scale unsupervised language models (LMs) acquire extensive world knowledge and certain reasoning capabilities, fine-grained regulation of their outputs remains a challenge because their training lacks any form of supervision. Common approaches to improve this controllability involve collecting human assessments of the relative merits of generated texts and subsequently adapting the unsupervised LM to reflect these preferences, often leveraging reinforcement learning from human feedback (RLHF). Yet, RLHF is a multi-stage and often unstable technique, beginning with the development of a reward model that approximates human choices, followed by reinforcement learning-based adaptation of the LM to maximize this proxy reward while constraining deviations from the pre-trained model. In this work, we propose a novel parameterization for the reward model within RLHF, facilitating direct derivation of the corresponding optimal policy in closed form. This reformulation allows us to address the classical RLHF problem with a straightforward classification objective. Our method, termed \textit{Direct Preference Optimization} (DPO), is robust, efficient, and lightweight, obviating the need for LM sampling during fine-tuning or intensive hyperparameter adjustment. Experimental results demonstrate that DPO enables LMs to adhere to human preferences at least as effectively as prior techniques. In particular, DPO surpasses PPO-based RLHF in moderating the sentiment of outputs, and performs on par with or better than established methods in tasks such as summarization and single-turn dialogue, all while being significantly simpler to deploy and train.
\end{abstract}

\section{Introduction}
Large-scale unsupervised language models (LMs) trained on extensive corpora have demonstrated a range of unexpected competencies~\citep{chowdhery2022palm, brown2020language, touvron2023llama,bubeck2023sparks}. These systems, however, learn from data created by humans whose intentions, expertise, and motivations span a broad spectrum. Not all such characteristics are desirable for imitation; for instance, while an AI-based code assistant ought to \textit{recognize} frequent programming errors in order to remediate them, it is preferable for the model’s code generation to favor the (sometimes infrequent) instances of expert-level coding present in its training distribution. Likewise, it is useful for a language model to \textit{know about} widespread misconceptions that half the population might hold, but we would not want it to endorse those inaccuracies in half of its responses to relevant prompts. Thus, carefully curating the model’s \emph{intended outputs and conduct} from among its expansive \textit{internalized knowledge and skill} is essential for developing AI agents that are reliable, high-performing, and amenable to control \citep{ouyang2022training}. Contemporary approaches predominantly employ reinforcement learning (RL) to shape LM behavior in accordance with human values, but as we shall demonstrate, the RL-oriented objective common to these frameworks can actually be minimized precisely using a straightforward binary cross-entropy loss, which streamlines the preference learning workflow.

In essence, current strategies imbue language models with preferred traits through tailored datasets reflecting human judgments about what is considered beneficial and safe. This process—often termed preference learning—follows an initial unsupervised pre-training phase leveraging a massive text corpus. Although supervised fine-tuning on curated demonstrations is a direct method for aligning models with human expectations, the dominant and most effective category of approaches is reinforcement learning from human (or artificial) feedback (RLHF/RLAIF; \citep{christiano2017deep,bai2022constitutional}). RLHF works by first constructing a reward model trained on human preference data, which is then used within an RL loop to guide the policy of a language model toward generating responses rated highly by the reward model, while constraining divergence from the pre-trained model. Despite yielding language models with state-of-the-art dialogue and programming performance, RLHF procedures are notably more involved than supervised learning alternatives: they necessitate training multiple LM instances and recurrently sampling policy outputs throughout the training cycle, leading to substantial additional computational expenditure.

In this work, we introduce a method to optimize language models in accordance with human preferences, eliminating the need for explicit reward modeling or reinforcement learning frameworks. We present \textit{{Direct Preference Optimization} (DPO)}, an approach that, while implicitly targeting the same optimization objective as standard RLHF algorithms (maximizing reward subject to a KL-divergence regularization), remains straightforward both in terms of implementation and training procedures. The core idea behind the DPO method is to adjust the model so as to increase the log odds of preferred responses compared to less preferred ones, while employing an adaptive, example-specific weighting mechanism that curbs the model collapse observed when using a basic probability ratio formulation. In line with prior work, DPO leverages a theoretical model of preferences (e.g., the Bradley-Terry model; \cite{bradley1952rankanalysis}) to gauge the concordance between a reward function and observed preference data. The distinction lies in how DPO utilizes this preference model: whereas established approaches first use it to define and minimize a loss over a reward model before optimizing a policy against the resulting learned reward, DPO directly expresses the preference loss in terms of the policy itself by reparameterizing the loss, thus bypassing the intermediary reward modeling step. As a result, given annotated human preference data for model outputs, DPO allows the policy to be learned via a binary cross entropy loss, efficiently training the model to maximize an implicit reward function that is consistent with the observed preferences.

The primary innovation of our work is Direct Preference Optimization (DPO), a simple, reinforcement learning–free approach for preference-based language model training. Empirically, we demonstrate that DPO matches or exceeds the effectiveness of existing techniques, such as PPO-based RLHF, across diverse preference-driven applications including sentiment modification, text summarization, and conversational tasks, for language models scaling up to 6B parameters.

\section{Related Work}

Self-supervised language models of greater size are capable of performing certain tasks in a zero-shot manner \citep{radford2019language} or through few-shot prompting strategies \citep{gpt3,megatron,chowdhery2022palm}. Nevertheless, enhancing their effectiveness on specific downstream tasks and improving their alignment with user instructions is often achieved by fine-tuning on curated collections of instructions and human-authored responses \citep{mishra-etal-2022-cross,sanh2022multitask,chung2022scaling,thoppilan2022lamda}. This technique, known as `instruction-tuning', facilitates broader generalization by LLMs to novel instructions beyond those seen in the fine-tuning corpus, and typically results in improved usability \citep{chung2022scaling}. While instruction tuning has proven successful, collecting \textit{relative} assessments of response quality from humans tends to be more straightforward than gathering expert-labeled demonstrations. As a result, subsequent research has explored fine-tuning LLMs using datasets comprising human preference judgments, thereby boosting model competency in areas such as translation \citep{kreutzer-etal-2018-reliability}, summarization \citep{stiennon2022learning,ziegler2020finetuning}, narrative generation \citep{ziegler2020finetuning}, and instruction adherence \citep{ouyang2022training,ramamurthy2023is}. These approaches generally begin by training a neural reward model to align with preference data, often leveraging frameworks such as the Bradley-Terry model \citep{bradley1952rankanalysis} for pairwise comparisons. Subsequent model fine-tuning is then performed to maximize the inferred reward signal, employing reinforcement learning algorithms like REINFORCE \citep{williams1992reinforce}, proximal policy optimization (PPO; \cite{schulman2017proximal}), or related approaches \citep{ramamurthy2023is}. Parallel developments make use of LLMs adapted for instruction-following via human feedback to generate additional artificial preference data focused on attributes such as safety or non-harmfulness \citep{bai2022constitutional}, relying on weak forms of human supervision such as rubric-based textual guidelines for annotation. These efforts represent a synthesis between two major research areas: training language models with reinforcement learning toward varied objectives~\citep{Ranzato2015SequenceLT,paulus2018a,wu2018learning}, and advancing general methodologies for learning from human preferences \citep{christiano2017deep,kupcsik2018learning}. In spite of the advantages offered by using relative human preferences, applying reinforcement learning for fine-tuning large language models introduces substantial practical difficulties; in response, this work proposes a theoretically sound method for optimizing with relative preferences that does not rely on RL.

The study of learning policies from preferences outside of linguistic applications has received attention within both bandit frameworks and reinforcement learning, resulting in a variety of proposed methodologies. In particular, the contextual dueling bandit (CDB; \cite{yue2012karmed,dudik2015contextual}) framework focuses on contextual bandit problems where feedback is provided via preferences or rankings over actions, as opposed to explicit rewards. In the absence of absolute reward signals, the theoretical treatment of CDBs introduces the concept of a \textit{von Neumann winner}—defined as a policy whose expected probability of prevailing over any alternative policy is at least 50\% \citep{dudik2015contextual}. A notable distinction is that preference annotations in CDBs are provided in an online manner, whereas typical practice in learning from human preferences involves using a static, offline batch of action pairs annotated with preferences \citep{yan2022human}. In a similar vein, \textit{preference-based RL} (PbRL) relies on binary preference information derived from an unknown underlying ‘scoring’ function instead of relying on conventional rewards \citep{BusaFekete2014,ruiz2023dueling}. Numerous algorithms have been developed for PbRL; some can leverage off-policy preference datasets, though these approaches generally entail first learning an explicit approximation to the latent scoring or reward function and subsequently performing policy optimization with respect to it \citep{jain2013learning,BusaFekete2014,christiano2017deep,sadigh2017active,kupcsik2018learning}. In contrast, we introduce a unified approach where policy optimization is carried out directly with the objective of satisfying observed preferences.

\section{Preliminaries}\label{section:prelims}

We provide an overview of the RLHF framework as described by \citeauthor{ziegler2020finetuning}, and later elaborated by works such as \citep{stiennon2022learning, bai2022training, ouyang2022training}. This process typically comprises three main stages: 1) supervised fine-tuning (SFT); 2) collection of preference data and reward function learning; and 3) reinforcement learning-based optimization.

\textbf{SFT}: The RLHF process commences with supervised fine-tuning, wherein a large pre-trained language model is further trained on a curated set of high-quality labeled examples targeting relevant downstream tasks (such as summarization or dialogue), yielding a model denoted as $\pi^\text{SFT}$.

\textbf{Reward Modelling Phase}: Next, the SFT-tuned model is provided with prompts $x$ to generate answer pairs $(y_1, y_2)\sim \pi^\text{SFT}(y \mid x)$. These answer pairs are subsequently evaluated by human annotators, who register a preference—denoted $y_w\succ y_l \mid x$—for one response over the other; here, $y_w$ and $y_l$ represent the chosen and non-chosen completions, respectively. The source of these preferences is presumed to be an inaccessible, latent reward model $r^*(y, x)$. Multiple approaches are available to model human preference data, with the Bradley-Terry (BT) \cite{bradley1952rankanalysis} model being a frequently used choice (and more general models, such as Plackett-Luce \citep{plackett1975analysis, luce2012individual}, can be employed if multiple ranked completions are available). According to the BT model, the probability distribution of observed preferences $p^*$ can be expressed as follows:

\begin{equation}\label{eq:bradley-terry}
    p^*(y_1\succ y_2 \mid x)=\frac{\exp\left(r^*(x, y_1)\right)}{\exp\left(r^*(x, y_1)\right) + \exp\left(r^*(x, y_2)\right)}.
\end{equation}

Given a fixed dataset of pairwise comparisons $\mathcal{D}=\bigl\{x^{(i)}, y_w^{(i)}, y_l^{(i)}\bigr\}_{i=1}^N$ drawn from the distribution $p^*$, we can introduce a reward model $r_{\phi}(x, y)$ with parameters $\phi$ and train it via maximum likelihood estimation. By re-casting the comparison task as a binary classification problem, we obtain the following negative log-likelihood loss:

\begin{equation}\label{eq:reward_model}
    \mathcal{L}_R(r_{\phi}, \mathcal{D}) = -\mathbb{E}_{(x, y_w, y_l)\sim \mathcal{D}}\bigl[\log \sigma(r_{\phi}(x, y_w)- r_{\phi}(x, y_l))\bigr]
\end{equation}

with $\sigma$ denoting the logistic sigmoid. In large language models, $r_{\phi}(x, y)$ is typically initialized with parameters from a supervised fine-tuned (SFT) model $\pi^\text{SFT}(y \mid x)$, augmented with an additional linear head applied atop the output of the final transformer layer to produce a single-valued reward estimate \cite{ziegler2020finetuning}. To reduce the variance in the learned reward, it is common practice to normalize so that $\mathbb{E}_{x,y\sim \mathcal{D}}\left[r_\phi(x, y)\right] = 0$ for each $x$.

\textbf{RL Fine-Tuning Phase}: In the subsequent reinforcement learning (RL) step, the trained reward model serves to provide evaluative feedback for the generation policy of the language model. Adhering to previous approaches~\citep{jaques2017sequence, jaques2020human}, the optimization target is expressed as

\begin{equation}\label{eq:RL}
\max_{\pi_{\theta}}  \mathbb{E}_{x\sim \mathcal{D}, y\sim \pi_{\theta}(y \mid x)}\bigl[r_{\phi}(x, y)\bigr] - \beta\mathbb{D}_{\textrm{KL}}\bigl[\pi_{\theta}(y\mid x)\mid \mid \pi_\text{ref}(y\mid x)\bigr],
\end{equation}

where $\beta$ determines how much the optimized policy $\pi_\theta$ may diverge from the reference policy $\pi_\text{ref}$ (usually the initial SFT model $\pi^\text{SFT}$). Typically, $\pi_\theta$ is initialized from $\pi^\text{SFT}$ as well. The KL penalty acts to limit policy shift from the region in which the reward model remains valid and ensures generative diversity, reducing risk of collapse to a small subset of high-reward completions. Owing to the inherently discrete outputs of sequence models, this objective cannot be directly differentiated and must be optimized using RL methods. The prevailing solution \citep{ziegler2020finetuning, stiennon2022learning, bai2022training, ouyang2022training} formulates the reward as ${r(x, y) = r_{\phi}(x, y) -\beta (\log \pi_{\theta}(y\mid x) - \log \pi_\text{ref}(y\mid x))}$, maximizing it with PPO \cite{schulman2017proximal}. 

\section{Direct Preference Optimization}\label{sec:DPO}

In light of the scalability and complexity challenges associated with RL algorithms for large-scale language model fine-tuning, we aim to develop a direct method for optimizing policies based on preferences. Distinct from conventional RLHF workflows that first estimate a reward model before applying RL, our method adopts a specific reward model parameterization which permits the direct computation of its optimal policy in closed-form, dispensing with the need for iterative RL updates. As detailed below, the central idea is to use an explicit analytic relationship linking reward functions to their optimal policies, allowing a transformation of the reward-model loss into a policy-centric loss. This variable substitution eliminates the need to fit an explicit, independent reward predictor, yet still ensures optimization according to established preference modeling frameworks, such

\textbf{Derivation of the DPO Objective.} We commence with the reinforcement learning objective defined in Eq.~\ref{eq:RL}, adopting a general reward function $r$, as has been standard in previous studies. Prior research~\citep{peters2007reinforcement, peng2019advantage, korbak2022reinforcement, go2023aligning} demonstrates that the optimizer of the reward maximization problem, subject to a KL-divergence constraint in Eq.~\ref{eq:RL}, admits the solution:

\begin{equation}\label{eq:op_policy}
    \pi_r(y\mid x) = \frac{1}{Z(x)}\pi_\text{ref}(y\mid x)\exp\left(\frac{1}{\beta}r(x, y)\right),
\end{equation}

where $Z(x) =\sum_{y}\pi_\text{ref}(y\mid x)\exp\left(\frac{1}{\beta}r(x, y)\right)$ serves as the partition function. For a complete derivation, refer to Appendix~\ref{app:derivation1}. Importantly, estimating the normalization constant $Z(x)$ remains challenging, even when substituting the reward $r$ with its maximum likelihood estimate $r_\phi$ for the true reward $r^*$, as highlighted by \citet{korbak2022reinforcement, go2023aligning}. Consequently, this parametric form is not always practical for direct use. However, Eq.~\ref{eq:op_policy} can be algebraically manipulated to represent the reward function as a function of its optimal policy $\pi_r$, the base policy $\pi_\text{ref}$, and the (unknown) $Z(x)$. Explicitly, by taking logarithms of both sides and rearranging terms, we find:

\begin{equation}\label{eq:main_eq}
    r(x,y) =\beta \log \frac{\pi_r(y\mid x)}{\pi_\text{ref}(y\mid x)} + \beta \log Z(x).
\end{equation}

This reparameterization can be employed for both the ground-truth reward $r^*$ and the optimal policy $\pi^*$. The Bradley-Terry preference model only depends on the reward differential between two outputs; that is, ${p^*(y_1 \succ y_2 \mid x) = \sigma(r^*(x, y_1) - r^*(x, y_2))}$. By inserting Eq.~\ref{eq:main_eq} for $r^*(x,y)$ into the preference formulation in Eq.~\ref{eq:bradley-terry}, the terms involving $Z(x)$ cancel, allowing the probability of human preference to be specified entirely in terms of $\pi^*$ and $\pi_\text{ref}$. As a result, the optimal policy $\pi^*$ that emerges from RLHF under the Bradley-Terry framework satisfies:

\begin{equation}\label{eq:objective}
    p^*(y_1\succ y_2 \mid x)=\frac{1}{1 + \exp\left(\beta \log \frac{\pi^*(y_2\mid x)}{\pi_\text{ref}(y_2\mid x)} - \beta \log \frac{\pi^*(y_1\mid x)}{\pi_\text{ref}(y_1\mid x)}\right)}
\end{equation}

The detailed derivation is provided in Appendix~\ref{app:derivation2}. Although this formulation is based on the Bradley-Terry model, a parallel analysis leads to corresponding expressions under the Plackett-Luce models~\citep{plackett1975analysis, luce2012individual}, which are elaborated in Appendix~\ref{app:plackett_luce_models}.

Given that the likelihood of observing a human preference is now framed in terms of the optimal policy (rather than the reward function), we may now define a maximum likelihood criterion for a parametrized policy $\pi_\theta$. In a manner analogous to reward-based modeling (see Eq.~\ref{eq:reward_model}), the resulting objective for policy optimization takes the form:

\begin{equation}\label{eq:optimum_model}
    \mathcal{L}_\text{DPO}(\pi_{\theta}; \pi_\text{ref}) = -\mathbb{E}_{(x, y_w, y_l)\sim \mathcal{D}}\left[\log \sigma \left(\beta \log \frac{\pi_{\theta}(y_w\mid x)}{\pi_\text{ref}(y_w\mid x

In this formulation, we construct an implicit reward through an alternative parameterization, such that the optimal policy corresponds directly to $\pi_\theta$. Additionally, as our method is mathematically equivalent to training a reparameterized Bradley-Terry model, it inherits desirable theoretical attributes, including consistency under appropriate assumptions on the distribution of preference data, as established in \cite{bong2022generalized}. We elaborate on these theoretical aspects of DPO and its connection to prior research in Section~\ref{sec:theory}.

\textbf{How does the DPO update operate?} To gain insight into the mechanics of DPO, it is instructive to study the gradient of its loss function, $\mathcal{L}_\text{DPO}$. The parameter gradient can be expressed as:

\begin{multline*}\label{eq:gradient}
    \nabla_\theta \mathcal{L}_\text{DPO}(\pi_\theta;\pi_\text{ref}) = \\ -\beta\mathbb{E}_{(x, y_w, y_l) \sim \mathcal{D}} \bigg[\underbrace{\sigma(\hat{r}_\theta(x, y_l) - \hat{r}_\theta (x, y_w))}_\text{greater emphasis when reward misranking occurs}\bigg[\underbrace{\nabla_\theta\log \pi(y_w \mid x)}_\text{promote $y_w$ occurrence} - \underbrace{\nabla_\theta\log\pi(y_l \mid x)}_\text{demote $y_l$ occurrence}\bigg]\bigg],
\end{multline*}

where the implicit reward $\hat{r}_\theta(x, y)$ is defined as $\beta \log \frac{\pi_\theta(y \mid x)}{\pi_\text{ref}(y \mid x)}$, leveraging both the trained language model $\pi_\theta$ and a reference model $\pi_\text{ref}$ (refer to Section~\ref{sec:theory} for further discussion). Conceptually, the gradient for $\mathcal{L}_\text{DPO}$ acts to upweight the probability assigned to preferred outputs $y_w$, while downweighting the probability of less-preferred responses $y_l$. Crucially, the contribution of each training sample is modulated by the extent to which the implicit reward model $\hat{r}_\theta$ overvalues the dispreferred completion, amplified by $\beta$, thus correcting more aggressively when the ordering of completions is violated—a behavior which also reflects the regularization imposed by the KL term. Empirical evidence from our experiments indicates that this adaptive weighting is vital, as omitting it (i.e., removing the weighting term) can result in degeneration of the language model, as shown in Appendix Table~\ref{tab:unlikelihood_generations}.

\textbf{DPO overview.} The standard DPO procedure consists of the following steps: 1) For each prompt $x$, generate completions $y_1, y_2 \sim \pi_\text{ref}(\cdot \mid x)$, and annotate these using human preferences to form an offline preference dataset $\mathcal{D} = \{x^{(i)}, y_w^{(i)}, y_l^{(i)}\}_{i=1}^N$; and 2) update the parameters of the language model $\pi_\theta$ by minimizing the DPO loss $\mathcal{L}_\text{DPO}$, utilizing the selected $\pi_\text{ref}$, dataset $\mathcal{D}$, and the specified $\beta$. In practical scenarios, it is often preferable to leverage publicly available preference datasets rather than collecting new samples and human annotations. Since these preference datasets are typically drawn from $\pi^\text{SFT}$, we initialize $\pi_\text{ref}$ with $\pi^\text{SFT}$ whenever this is feasible. If $\pi^\text{SFT}$ cannot be obtained, $\pi_\text{ref}$ is instead initialized by maximizing the likelihood over the preferred completions ${(x, y_w)}$, i.e., setting ${\pi_\text{ref} = \argmax_{\pi}\mathbb{E}_{x, y_w \sim \mathcal{D}}\left[\log \pi(y_w \mid x)\right]}$. This approach serves to minimize the distributional gap between the inaccessible true reference distribution and the reference model $\pi_\text{ref}$ used within DPO. Additional implementation details and choices of hyperparameters are discussed in Appendix~\ref{app:implementation}.

\section{Theoretical Analysis of DPO} In this section, we provide a deeper examination of the DPO framework, furnishing theoretical justification and drawing connections between the advantages of DPO and the shortcomings of actor-critic methods commonly employed in RLHF settings (for instance, PPO~\cite{schulman2017proximal}).

\label{sec:theory}

\subsection{Your Language Model Is Secretly a Reward Model} DPO eliminates the need for explicit reward model fitting and subsequent RL-based policy optimization by consolidating both steps into a unified maximum likelihood training procedure. Specifically, the objective in Eq. \ref{eq:main_eq} mirrors the Bradley-Terry model with a reward defined as $r^*(x, y) = \beta \log\frac{\pi^*_\theta(y \mid x)}{\pi_\text{ref}(y \mid x)}$, so that optimizing the parametric model $\pi_{\theta}$ directly corresponds to reward model training in Eq. \ref{eq:reward_model} under a suitable variable transformation. In what follows, we develop the theoretical foundation for this reparameterization, demonstrate that it maintains the full expressive power of the reward model class, and prove that it admits exact recovery of the optimal policy. To begin, we introduce an equivalence relation among reward functions.

\begin{definition}
We define two reward functions $r(x, y)$ and $r'(x, y)$ as equivalent if and only if ${r(x, y)-r'(x, y) = f(x)}$ for some function $f$.     
\end{definition}

It is straightforward to verify that this relation is an equivalence relation, leading to a partition of the space of reward functions into equivalence classes. From this, we can assert the following lemmas:

\begin{lemma}\label{lemma:same_prefrence} Within the context of the Plackett-Luce model, including the Bradley-Terry formulation, any two reward functions belonging to the same equivalence class will induce identical distributions over preferences.
\end{lemma}

\begin{lemma}\label{lemma:same_policy}
    Within the constrained RL setting, reward functions drawn from a common equivalence class result in the same optimal policy.
\end{lemma}

The arguments are elementary and omitted here, with full proofs deferred to Appendix \ref{app:lemma1}. The first lemma captures a classic identifiability problem associated with the Plackett-Luce family of models \cite{plackett1975analysis}. Because this form of under-specification arises, it becomes necessary to introduce extra constraints in order to guarantee any identifiability for MLE solutions to Eq. \ref{eq:reward_model} \cite{bong2022generalized}. The second lemma implies that any member of a given equivalence class produces the same optimal policy in constrained RL, hence we only need to recover any single representative reward function from the optimal equivalence class. The following Theorem, proved in Appendix~\ref{app:thm1}, formalizes this insight:

\begin{theorem}\label{thm:main}
    Provided mild regularity conditions, all reward equivalence classes consistent with the Plackett-Luce (and specifically the Bradley-Terry) framework admit a representation of the form ${r(x, y) = \beta \log \frac{\pi(y\mid x)}{\pi_\text{ref}(y\mid x)}}$, where $\pi(y\mid x)$ is a model and $\pi_\text{ref}(y \mid x)$ is a fixed reference model.
\end{theorem}

\begin{sproof}
    Let $r(x, y)$ be an arbitrary reward function, which induces the optimal model $\pi_r(y \mid x)$ as defined in Eq. \ref{eq:op_policy}. We aim to demonstrate that there exists a representative in the equivalence class of $r$ that admits the desired reparameterized form. To do so, consider the projection mapping $f$ defined as  
\begin{equation}
    f(r; \pi_\text{ref}, \beta)(x, y) = r(x, y) - \beta\log\sum_{y}\pi_\text{ref}(y\mid x)\exp\left(\frac{1}{\beta}r(x, y)\right)
\end{equation}
This operator $f$ normalizes the reward by subtracting the log partition function (computed under $\pi_r$), and as this normalization term depends only on $x$, $f(r; \pi_\text{ref}, \beta)(x, y)$ lies in the same equivalence class as $r(x, y)$. Substituting $r$ in terms of the right-hand side of Eq.~\ref{eq:main_eq}, valid for any reward, yields $f(r; \pi_\text{ref}, \beta)(x, y) = \beta \log \frac{\pi_r(y\mid x)}{\pi_\text{ref}(y\mid x)}$. Thus, this construction produces an explicit element in the equivalence class with the reparameterized structure, demonstrating that the restriction imposed by this representation incurs no loss of generality.
\end{sproof}

Theorem~\ref{thm:main} can also be interpreted as pinpointing the precise reward function chosen by the DPO reparameterization within each equivalence class. Specifically, it selects the reward function that fulfills:

\begin{equation}\label{eq:lag_p}
     \sum_{y}\underbrace{\pi_\text{ref}(y\mid x)\exp\left(\frac{1}{\beta}r(x, y)\right)}_{=\pi(y\mid x)\text{, using Thm.~\ref{thm:main} reparam.}} = 1,
\end{equation}

which ensures that $\pi(y\mid x)$ constitutes a proper probability distribution—non-negative and summing to one. Furthermore, by referencing Eq.~\ref{eq:op_policy}, it becomes apparent that Eq.~\ref{eq:lag_p} essentially defines the partition function for the optimal policy corresponding to the reward function $r(x, y)$. The central idea underpinning the DPO algorithm is to enforce suitable constraints within the under-identified Plackett-Luce (notably including the Bradley-Terry) preference model family. This approach retains the same expressive power for representable reward models, while simultaneously guaranteeing that the optimal policy described in Eq. \ref{eq:op_policy} can always be computed explicitly for any prompt $x$.

\subsection{Instability of Actor-Critic Algorithms}  
Our framework further enables the identification of sources of instability within typical actor-critic methods utilized for RLHF, like PPO. Adhering to the RLHF setup, we concentrate on the RL fine-tuning phase outlined in Section \ref{section:prelims}. We draw parallels to the control-as-inference paradigm \cite{levine2018reinforcement} for the constrained RL scenario set forth in \ref{eq:RL}. In this setting, with a parameterized policy $\pi_{\theta}(y\mid x)$, the objective is to minimize $\mathbb{D}_{\text{KL}}[\pi_{\theta}(y|x) \mid \mid \pi^*(y\mid x)]$, where $\pi^*$ is the reward-induced optimal policy given by Eq.~\ref{eq:optimum_model} using $r_{\phi}(y, x)$. With further manipulation, the resulting optimization problem takes the following form:

\begin{equation}\label{eq:AC}
    \max_{\pi_{\theta}}\mathbb{E}_{\pi_{\theta}(y\mid x)}\bigg[\underbrace{r_{\phi}(x, y) -\beta\log\sum_{y}\pi_\text{ref}(y\mid x)\exp\left(\frac{1}{\beta}r_{\phi}(x, y)\right)}_{f(r_{\phi}, \pi_\text{ref}, \beta)} - \underbrace{\beta\log\frac{\pi_{\theta}(y\mid x)}{\pi_\text{ref}(y\mid x)}}_{\text{KL}}\bigg]
\end{equation}

The objective considered here aligns with that optimized in previous research 
\citep{ziegler2020finetuning, stiennon2022learning, bai2022training, ouyang2022training}, where the DPO-equivalent reward is applied to the reward class $r_{\phi}$. In this framework, the normalization term in $f(r_{\phi}, \pi_\text{ref}, \beta)$ can be viewed as the soft value function corresponding to the reference policy $\pi_\text{ref}$. Although this normalization term does not influence the final optimal solution, omitting it can introduce substantial variance into the policy gradient, potentially resulting in unstable training. One method to address the normalization term is to approximate it with a learned value function, though this approach can pose significant optimization challenges. Historically, other works have normalized the rewards by employing a human completion baseline, functioning as a single-sample Monte Carlo approximation of the normalizer. By contrast, the DPO reparameterization provides a reward function inherently independent of such baselines.

\section{Experiments}
Here, we undertake empirical analyses to assess the effectiveness of DPO in directly training policies from preference data. We start with a tightly controlled text-generation task to investigate: how effectively does DPO balance reward maximization and KL-divergence minimization with respect to the reference policy, when contrasted with standard preference learning techniques like PPO? Subsequently, we extend our evaluation to encompass larger model architectures and more challenging RLHF scenarios, such as summarization and dialogue. The results indicate that, with minimal hyperparameter adjustment, DPO frequently matches or outperforms strong methods—including RLHF via PPO and selection of the top trajectory from $N$ samples under a learned reward model. Before detailing these findings, we outline the experimental methodology, with further information provided in Appendix~\ref{app:exp_details}.

\textbf{Tasks.} We investigate three distinct open-ended text generation tasks in our experimental setup. In all cases, learning algorithms derive a policy using a preference dataset $\mathcal{D}=\bigl\{x^{(i)}, y_w^{(i)}, y_l^{(i)}\bigr\}_{i=1}^N$. For \textbf{controlled sentiment generation}, $x$ serves as an initial segment of a movie review extracted from the IMDb dataset \cite{maas-EtAl:2011:ACL-HLT2011}, and the objective for the policy is to produce $y$ exhibiting positive sentiment. To facilitate controlled measurement, we synthetically generate preference pairs for model generations using a pre-trained sentiment classifier, ensuring that $p(\text{positive}\mid x,y_w)>p(\text{positive}\mid x,y_l)$. For SFT, GPT-2-large is fine-tuned until performance plateaus using training set reviews from IMDB (additional methodological specifics are found in App~\ref{app:sentiment_details}). In the \textbf{summarization} task, $x$ is a Reddit forum post, with the policy responsible for producing a summary $y$ that captures the principal points. Adhering to established methodologies, we utilize the Reddit TL;DR dataset \citep{volske-etal-2017-tl} along with a set of human preferences curated by \citeauthor{stiennon2022learning}. The SFT model is trained on forum post summaries authored by humans\footnote{\url{https://huggingface.co/CarperAI/openai_summarize_tldr_sft}}, employing the TRLX framework \citep{leandro_von_werra_2023_7790115} for RLHF training. Note that the human-labeled preference data by \citeauthor{stiennon2022learning} was created from outputs of another SFT model trained in a similar fashion. The \textbf{single-turn dialogue} task involves $x$ as a user-issued query, spanning topics from scientific inquiries to personal advice. The policy's aim is to generate a response $y$ that is informative and engaging; for this we rely on the Anthropic Helpful and Harmless dialogue dataset \citep{bai2022training}, which comprises 170,000 dialogues between people and automated assistants. Each conversation concludes with two possible model-generated responses, one marked as preferred by a human annotator. Since a pre-trained SFT model is not provided in this context, we develop an SFT model by fine-tuning a generic language model solely on the human-preferred responses.

\textbf{Evaluation.} We employ two distinct evaluation strategies in our experimental framework. To systematically assess how each algorithm optimizes the constrained reward maximization objective, we analyze their trade-off between attained reward and KL-divergence from the reference policy in the controlled sentiment generation scenario; here, the entire reward–divergence frontier is accessible since we utilize a ground-truth reward function (namely, a sentiment classifier). By contrast, practical scenarios lack explicit access to the true reward function. As a result, our main evaluation in such real-world settings consists of measuring the \textit{win rate} of each algorithm when pitted against a baseline policy, leveraging GPT-4 to approximate human preferences in both summarization and single-turn dialogue contexts—measuring summary quality and helpfulness of responses, respectively. For summarization, the baseline comprises reference summaries from the test split; for dialogue, it is the annotator-preferred answer from the evaluation dataset. Although some literature suggests language models surpass traditional metrics as automatic evaluators \citep{Chen2023ExploringTU}, we additionally validate our use of GPT-4 by conducting a human study, as detailed in Sec.~\ref{sec:human-judgments}. The findings show a strong alignment between GPT-4 and human judgments, with human-GPT-4 agreement rates generally on par with, or exceeding, inter-annotator reliability.

\textbf{Methods.} Beyond DPO, our empirical study compares a variety of established approaches for aligning language models to human preferences. The simplest baselines include zero-shot prompting with \textbf{GPT-J} \citep{gpt-j} in summarization and 2-shot prompting with \textbf{Pythia-2.8B} \citep{biderman2023pythia} for dialogue. We further evaluate models trained by supervised fine-tuning (\textbf{SFT}) and the \textbf{Preferred-FT} method, which involves supervised learning using the human-preferred response $y_w$—originating from SFT for sentiment and summarization, or from a generic language model for dialogue tasks. Another comparative model is based on \textbf{Unlikelihood}~\citep{welleck2019neural} training, where the objective maximizes the likelihood of $y_w$ while explicitly reducing the probability of $y_l$, weighted by an optional `unlikelihood' coefficient $\alpha\in[0,1]$. Furthermore, we incorporate \textbf{PPO} \citep{schulman2017proximal}, utilizing a reward model trained on preference annotations, and also an oracle variant, \textbf{PPO-GT}, that benefits from access to ground-truth rewards in the sentiment-controlled setting. Within sentiment experiments, we explore both an off-the-shelf PPO-GT implementation \cite{leandro_von_werra_2023_7790115} and a customized variant featuring reward normalization and additional hyperparameter tuning—optimizations we also apply to the conventional PPO method when training with learned rewards. Lastly, we include the \textbf{Best of $N$} heuristic, which samples $N$ completions from the SFT model (or Preferred-FT in the case of dialogue) and selects the top response based on a reward model trained from preference data. While this method may yield strong empirical results by decoupling reward modeling from PPO’s optimization dynamics, its utility is constrained by its high computational requirements, given that it necessitates generating $N$ completions per test example.

\subsection{How well can DPO optimize the RLHF objective?}

In standard RLHF approaches, the reward maximization objective is regularized by a KL constraint, ensuring the policy optimizes for high reward without straying too far from the baseline policy. Thus, when assessing different methods, it is important to jointly consider both the level of reward attained and the extent of the KL divergence; a marginally higher reward at the cost of a much greater KL may not be preferable. Figure~\ref{fig:frontier-tldr-main} illustrates the tradeoff between reward and KL for a range of methods applied to the sentiment task. For each approach, we conduct several training runs, each corresponding to a different setting of the conservativeness hyperparameter (target KL $\in\{3,6,9,12\}$ for PPO, $\beta \in \{0.05,0.1,1,5\}$, $\alpha\in\{0.05,0.1,0.5,1\}$ for unlikelihood, random seeds for preferred-FT). In total, this hyperparameter sweep results in 22 training runs. At intervals of every 100 training steps until convergence, we evaluate all trained policies on a set of test prompts, measuring the mean reward according to the true reward model and computing the average KL at the sequence level\footnote{Namely, summing KL divergences over all timesteps in a sequence.} relative to the reference policy, $\text{KL}\left(\pi\mid \mid \pi_\text{ref}\right)$. The results demonstrate that DPO achieves the most advantageous reward-KL tradeoff, with superior rewards at relatively low KL values. This finding is significant for several reasons. Most notably, although both DPO and PPO are designed to optimize the same objective, DPO is markedly more effective: its reward-versus-KL curve consistently surpasses that of PPO. Furthermore, DPO yields a better frontier compared to PPO, \emph{even when PPO is given access to the ground-truth reward signal} (PPO-GT).

\subsection{Can DPO scale to real preference datasets?}
\label{sec:dpo-real-datasets}
We further investigate the effectiveness of DPO in the context of fine-tuning for summarization and single-turn dialogue tasks. In summarization, traditional automatic metrics like ROUGE often do not align well with human judgments~\citep{stiennon2022learning}. Earlier studies have demonstrated that language models fine-tuned with PPO using human preference data can yield superior summaries. Our evaluation involves generating outputs on the TL;DR summarization dataset’s test partition, then measuring the average win rate of these outputs compared to reference completions. For each method, outputs are sampled with temperatures ranging between 0.0 and 1.0, and the resulting win rates are presented in Figure~\ref{fig:frontier-tldr-main} (right). The fine-tuning process for DPO, PPO, and Preferred-FT starts from the same GPT-J SFT model\footnote{\url{https://huggingface.co/CarperAI/openai_summarize_tldr_sft}}. At a sampling temperature of 0.0, DPO achieves a win rate of about 61\%, outperforming PPO, which reaches around 57\% at its optimal temperature. DPO also surpasses the best-performing $N$-shot baseline in terms of peak win rate. It is important to point out that DPO’s $\beta$ parameter was not exhaustively tuned, which suggests that its performance could be further improved. Additionally, DPO displays significantly greater robustness to changes in sampling temperature relative to PPO, whose win rate drops to that of the unmodified GPT-J model as temperature increases. Preferred-FT yields only marginal gains over the original SFT model. In Section~\ref{sec:human-judgments}, we conduct direct human preference comparisons between DPO and PPO; DPO samples generated with a temperature of 0.25 were chosen over PPO samples at temperature 0.0 in 58\% of cases.

For single-turn dialogues, our assessment focuses on the subset of the Anthropic HH test set \citep{bai2022training} that includes only a single human-assistant exchange. To evaluate performance, GPT-4 assessments reference the test set's preferred completions when determining win rates across the various approaches. In the absence of an established SFT model for this context, our workflow begins with a pre-trained Pythia-2.8B model. We first train a reference model via Preferred-FT, using the selected completions to ensure the generated outputs align with the model's output distribution, followed by further training utilizing DPO. Comparisons are also made with the strongest out of 128 Preferred-FT completions (as our analyses indicate that the Best of $N$ approach levels off at $N=128$ for this evaluation; see Appendix Figure~\ref{fig:best-of-n}) and with a 2-shot prompted version of the base Pythia-2.8B model. Our findings demonstrate that DPO matches or outperforms the best temperature settings across these methods. 

Additionally, we assess an RLHF model optimized with PPO on the Anthropic HH data \footnote{\url{https://huggingface.co/reciprocate/ppo_hh_pythia-6B}}, as provided by a reputable source \footnote{\url{https://github.com/CarperAI/trlx/tree/main/examples/hh}}, but are unable to identify any prompt or sampling temperature yielding superior results compared to the unmodified Pythia-2.8B. Drawing on results from TL;DR and acknowledging that both PPO and Best of 128 maximize the same reward objective, we regard Best of 128 as an approximate surrogate for PPO's performance. In summary, DPO stands out as the only efficient technique to surpass the preferred completions from the Anthropic HH test split, delivering outcomes on par with or better than the resource-intensive Best of 128 method. Furthermore, Figure~\ref{fig:dialogue-main} highlights DPO’s rapid convergence to optimal performance.

\subsection{Generalization to a new input distribution}

To more rigorously assess how PPO and DPO handle changes in input distributions, we analyze their performance on data from outside the training domain. Specifically, we test the policies learned from our Reddit TL;DR summarization study on the CNN/DailyMail dataset's test set, consisting of news articles \citep{nallapati-etal-2016-abstractive}. This evaluation utilizes the optimal sampling temperatures identified during TL;DR experiments (0 and 0.25), and the findings are reported in Table~\ref{tab:ood}. For evaluation, we measure GPT-4's win rate versus ground-truth summaries, maintaining the same GPT-4 (C) prompt as in the Reddit TL;DR study, but substituting “forum post” with “news article.” On this out-of-distribution task, DPO maintains a clear advantage over the PPO policy. These observations offer preliminary support that DPO is capable of generalizing to new domains as effectively as PPO, despite DPO not leveraging the unlabeled Reddit TL;DR prompts that benefit PPO.

\subsection{Validating GPT-4 judgments with human judgments}
\label{sec:human-judgments}
To confirm the robustness of GPT-4’s evaluation, we carry out a human subject experiment based on results from the TL;DR summarization task and two distinct prompts for GPT-4. The \textbf{GPT-4 (S)} (simple) prompt requests a judgment on which summary better captures the main information in a post. The \textbf{GPT-4 (C)} (concise) prompt, in contrast, also includes conciseness as a criterion, which addresses our observation that GPT-4, when given the (S) prompt, tends to favor overly lengthy and repetitive summaries compared to human preferences. The full text of each prompt is given in Appendix~\ref{app:prompts}. We select three methods for head-to-head comparisons against PPO with its optimal (greedy) temperature: the top-performing (DPO, temp. 0.25), the lowest-performing (PPO, temp. 1.0), and a mid-range model (SFT, temp. 0.25), aiming for a broad range of summary qualities. For these method pairs, and focusing especially on DPO and PPO-1, we obtain multiple independent human judgments (where feasible, given the available raters). Our results indicate that, for both prompts, the degree of agreement between GPT-4 and human raters matches inter-rater agreement among humans themselves, implying that GPT-4 can serve as an adequate stand-in for human raters. Of note, the \textbf{GPT-4 (C)} prompt yields results that most closely reflect human preferences, leading us to rely on this prompt for the main analyses in Section~\ref{sec:dpo-real-datasets}. Further information, such as details of the rating interface and a roster of human evaluators, can be found in Appendix~\ref{app:human-study}.

\section{Discussion}
Learning from preferences provides a robust and efficient approach for developing capable and aligned language models. In this work, we proposed DPO, a streamlined training framework that leverages preference data for training language models without relying on reinforcement learning. Rather than adapting preference learning to fit within a typical RL framework and utilizing standard RL algorithms, DPO establishes a correspondence between language model policies and reward structures, allowing the direct optimization of a language model toward human preferences via a straightforward cross-entropy objective—eschewing the need for reinforcement learning and avoiding loss of expressiveness. Remarkably, with almost no hyperparameter adjustment, DPO achieves comparable or improved results relative to prevalent RLHF methods, such as those based on PPO. As a result, DPO significantly lowers the entry barriers for training language models on human feedback.

\textbf{Limitations \& Future Work.} This study opens up several avenues for future investigation. One open question is how well the DPO-trained policy generalizes to unseen data in comparison to models trained with explicit reward signals. Preliminary evidence indicates DPO policies offer generalization on par with PPO-trained counterparts, yet a more exhaustive analysis remains to be conducted. Another promising area is examining whether DPO’s self-labeling capability can facilitate the use of prompts lacking labels. Additionally, further work is needed to understand manifestations of reward over-optimization within the direct preference optimization framework; for instance, the modest drop in performance observed in Figure~\ref{fig:dialogue-main}-right may be symptomatic of this. While our experiments extend to models with up to 6B parameters, scaling DPO to substantially larger, state-of-the-art architectures warrants further exploration. On the topic of evaluation, we noticed that win rates produced by GPT-4 can vary depending on prompt formulation, suggesting that developing strategies for obtaining consistent and reliable automated evaluations is an important next step. Beyond the domain of language modeling, DPO presents a range of potential applications, such as optimizing generative models in various modalities through preference-based training.